% =========================================================
% Proof file: prf-zero-scalar-times-vector.tex
% =========================================================

\newpage
\phantomsection
\label{prf:zero-scalar-times-vector}

\begin{remark*}[Return]
\hyperref[thm:zero-scalar-times-vector]{Return to Theorem}
\end{remark*}

\begin{theorem*}[$0v=0$]
For every $v \in V$,
\[
0v=0.
\]
\end{theorem*}

\begin{proof}
\textbf{Professional Standard Proof.}~\\
Let $v \in V$. By distributivity,
\[
(0+0)v=0v+0v.
\]
Because $0+0=0$ in $\mathbf{F}$, this becomes
\[
0v=0v+0v.
\]
Add the additive inverse of $0v$ to both sides:
\[
-(0v)+0v=-(0v)+(0v+0v).
\]
By associativity,
\[
0=\bigl(-(0v)+0v\bigr)+0v=0+0v=0v.
\]
Therefore
\[
0v=0.
\]
\end{proof}

\begin{proof}
\textbf{Detailed Learning Proof.}~\\

\textbf{Step 1.} Fix an arbitrary vector.\\
Let $v \in V$. The goal is to determine the vector $0v$.

\textbf{Step 2.} Use distributivity over scalar addition.\\
The vector space axiom gives
\[
(a+b)v=av+bv
\]
for all $a,b \in \mathbf{F}$. Apply this with $a=0$ and $b=0$:
\[
(0+0)v=0v+0v.
\]

\textbf{Step 3.} Simplify the scalar expression.\\
In the field $\mathbf{F}$,
\[
0+0=0.
\]
Hence the previous equation becomes
\[
0v=0v+0v.
\]

\textbf{Step 4.} Isolate $0v$ by adding its additive inverse to both sides.\\
Because $0v \in V$, it has an additive inverse $-(0v)$. Add $-(0v)$ to both sides:
\[
-(0v)+0v=-(0v)+(0v+0v).
\]

\textbf{Step 5.} Simplify both sides.\\
On the left,
\[
-(0v)+0v=0.
\]
On the right, associativity gives
\[
-(0v)+(0v+0v)=\bigl(-(0v)+0v\bigr)+0v=0+0v=0v.
\]
Therefore
\[
0=0v.
\]

\textbf{Step 6.} State the conclusion.\\
Thus scalar multiplication by the scalar zero yields the zero vector:
\[
0v=0.
\]
\end{proof}

\begin{remark*}[Proof structure]
The proof is direct. Distributivity first produces an equation of the form $w=w+w$ with $w=0v$, and the additive inverse of $w$ is then used to force $w=0$.
\end{remark*}

\begin{remark*}[Dependencies]
\hfill
\begin{itemize}
  \item \hyperref[def:vector-space]{Definition of Vector Space}
  \item \hyperref[thm:unique-additive-inverse]{Unique Additive Inverse}
\end{itemize}
\end{remark*}